% ============================================================================
% Libraries: what is used from each one and why.
% ============================================================================
\section{Bibliotecas utilizadas}\label{sec:librerias}

\Cref{tab:librerias} resume qué se usa de cada biblioteca. La regla general es usar funciones
compiladas (OpenCV, SciPy, NumPy vectorizado) para todo lo que se repite por píxel o por cuadro, y
dejar en Python solo la lógica de decisión.

\begin{table}[H]
  \centering
  \caption{Bibliotecas, funciones usadas y para qué. Versiones probadas entre paréntesis.}
  \label{tab:librerias}
  \small
  \begin{tabular}{@{}p{2.4cm}p{5.6cm}p{6.6cm}@{}}
    \toprule
    Biblioteca & Funciones & Uso \\
    \midrule
    \textbf{OpenCV} \newline (\texttt{cv2} 5.0) &
      \texttt{VideoCapture}, \texttt{CAP\_PROP\_*} &
      abrir el video, leer metadatos (cuadros, fps, tamaño) y saltar a un cuadro \\
    & \texttt{VideoWriter}, \texttt{VideoWriter\_fourcc} & escribir videos recortados o anotados \\
    & \texttt{cvtColor}, \texttt{resize} (\texttt{INTER\_AREA}) & pasar a grises y reducir la imagen a la escala de análisis \\
    & \texttt{filter2D} & convolución con los núcleos de disco y anillo (filtro de anillo, \cref{sec:deteccion}) \\
    & \texttt{warpPolar} & imagen en coordenadas polares alrededor de cada centro: cobertura del anillo y firma de rotación \\
    & \texttt{GaussianBlur}, \texttt{threshold}, \texttt{Canny}, \texttt{convertScaleAbs}, \texttt{dilate} & operaciones de la pestaña \emph{Procesar} y calibración \\
    & \texttt{circle}, \texttt{line}, \texttt{arrowedLine}, \texttt{putText}, \texttt{ellipse} & dibujo de la superposición y del video anotado \\
    \midrule
    \textbf{trackpy} \newline (0.7) &
      \texttt{locate} & máximos locales del mapa de respuesta con centroide subpíxel (\cref{sec:deteccion}) \\
    & \texttt{link}, \texttt{filter\_stubs} & vinculación simple de trayectorias en la pestaña \emph{Procesar} (no se usa en el seguimiento con $N$ fijo) \\
    \midrule
    \textbf{SciPy} \newline (1.18) &
      \texttt{optimize.linear\_sum\_assignment} & asignación óptima candidato--robot (algoritmo húngaro) \\
    & \texttt{spatial.cKDTree} & búsqueda de vecinos para confirmar robots nuevos por persistencia \\
    & \texttt{signal.savgol\_filter} & derivadas suavizadas: velocidad y $\omega$ \\
    & \texttt{sparse.coo\_matrix}, \texttt{sparse.linalg.spsolve} & mínimos cuadrados del ángulo acumulado \\
    \midrule
    \textbf{NumPy} \newline (2.5) &
      \texttt{fft.fft}, \texttt{fft.ifft} & armónicos de la firma de rotación y correlación circular \\
    & arreglos, \texttt{interp}, \texttt{savez\_compressed}, \texttt{load} & todo el cálculo numérico, interpolación de huecos y sesiones \texttt{.npz} \\
    \midrule
    \textbf{pandas} \newline (3.0) & \texttt{DataFrame}, \texttt{pivot}, \texttt{to\_csv}, \texttt{read\_csv} & tablas de candidatos y de resultados, formato ancho, CSV \\
    \midrule
    \textbf{PySide6} \newline (Qt 6.11) & \texttt{QMainWindow}, \texttt{QTabWidget}, widgets de formulario &
      interfaz \\
    & \texttt{QThread} + señales & tareas largas sin congelar la ventana \\
    & \texttt{QTimer} & reproducción del video \\
    & \texttt{QImage}, \texttt{QPainter} & visor con zoom y edición del rectángulo \\
    \midrule
    \textbf{Matplotlib} & \texttt{imshow}, \texttt{histogram2d} (NumPy) & solo en \texttt{comparar\_ensayos.py}: mapas de ocupación \\
    \bottomrule
  \end{tabular}
\end{table}

\paragraph{Por qué trackpy solo para localizar.} \texttt{trackpy.locate} es robusto para encontrar
máximos y refinarlos a precisión subpíxel, y por eso se usa. Su vinculador (\texttt{link}) no sabe que
la cantidad de robots es fija: ante un robot que desaparece unos cuadros crea una trayectoria nueva,
y ante un falso positivo, también. Por eso el seguimiento del ensayo usa un algoritmo propio con
$N$ fijo (\cref{sec:seguimiento}), y \texttt{link} queda solo como opción rápida en \emph{Procesar}.

\paragraph{Por qué OpenCV para el video.} \texttt{cv2.VideoCapture} usa FFmpeg internamente: lee
cualquier formato habitual y permite acceso aleatorio por número de cuadro, que necesitan el
reproductor, la calibración (cuadros salteados) y la revisión de cuadros con problemas.
